So much for \textsc{Kritter}'s opinion, or rather for the part of it directly connected with the subject of my present criticism. I shall return to the second part of the opinion, concerning the increase in expenditure arising from orphans' pensions, in the second division of this memorandum. That the first part of \textsc{Kritter}'s opinion touches only the surface of the subject will, I think, have been sufficiently shown by the preceding discussion. Reflections or conjectures about how \textsc{Kritter} came to give so superficial an opinion, not even in agreement with his other public statements, would take me too far here. Nevertheless, I am inclined to believe that if, instead of restricting matters to putting two specific questions,
\begin{quote}
one of which, moreover, admitted no definite answer at all in the sense intended by the questioner, they had allowed p. \textsc{Kritter} to take part in the further deliberations, in whatever form, or had at least consulted him on the suitability of the plan they thought could be based on his opinion, with a full account of all the circumstances,
\end{quote}
\noindent the question of the rule for all future pension increases would not have been so hastily snapped over the knee as it was\footnote{The number of instances of lax handling of business cited or hinted at in my account above could easily have been increased further, but I considered this as unnecessary as it was unpleasant.}. But none of this was done. \textsc{Kritter} received a fee of 4 Thl. 24 mgr. from the widows' fund for his opinion, just as P. received, in the ministerial rescript of 20 November 1794, a commendation for his well-prepared memorandum.

This rescript approved the proposals made, but stipulated, concerning the increase in pensions, that if, contrary to expectation, \textit{unfortunate circumstances} should later make a reduction in pensions necessary, the widows then present must accept it. Ratification on each occasion was also reserved for future pension increases to be made under the progression rule. \textsuperscript{.} It seems noteworthy that only a necessity arising from unfortunate circumstances is considered here, and not one that might arise from too large a number of widows. Was such a case considered impossible, or was the regulation understood in the same sense as \textsc{Sch.} understood it, with a reduction of the pension in such a case regarded as self-evident? The same reservation was repeated when the first increase was ratified (1799 May 23), but only circumstances that might make another reduction necessary were mentioned in general, without the qualification ‘unfortunate’. In later instances of ratification, so far as I have been able to find, the reservation was not repeated.

